\chapter{Tecnologie utilizzate}
\label{cap:tecnologie}

VisionCaption combina tecnologie web, computer vision e intelligenza
artificiale. In questo capitolo vengono presentate la natura e le
caratteristiche principali delle tecnologie adottate. Il modo in cui esse sono
integrate nell'architettura e nelle due modalità operative viene invece
affrontato nei capitoli di approfondimento e di sviluppo del sistema.

\section{Linguaggio e gestione del progetto}

\subsection{Python}
Python è un linguaggio di programmazione ad alto livello, multipiattaforma e di
uso generale. Utilizza una tipizzazione dinamica e dispone di costrutti per la
programmazione procedurale, a oggetti e funzionale. Il suo ecosistema comprende
numerose librerie per lo sviluppo web, il calcolo numerico, la computer vision
e il machine learning.

Python supporta inoltre la programmazione asincrona attraverso il modello
\texttt{async}/\texttt{await}, che consente di coordinare operazioni di
ingresso e uscita senza associare necessariamente un thread a ciascuna
operazione.

\subsection{uv}
\textit{uv} è uno strumento per la gestione di progetti e pacchetti Python.
Permette di creare ambienti virtuali, risolvere e installare le dipendenze e
generare un file di lock. Quest'ultimo registra le versioni esatte dei
pacchetti, rendendo riproducibile l'ambiente di esecuzione.

\subsection{Hatchling}
Hatchling è un backend di build conforme allo standard PEP 517. Il suo compito
è trasformare il codice sorgente e i metadati di un progetto Python in
distribuzioni installabili, come wheel e source distribution.

\section{Tecnologie web e comunicazione}

\subsection{Browser web}
Il browser è l'ambiente software che interpreta ed esegue le applicazioni web.
Oltre a HTML, CSS e JavaScript, espone API standard per accedere alle
funzionalità del dispositivo, tra cui fotocamera, riproduzione audio,
connessioni di rete e disegno bidimensionale. L'impiego di standard web
permette di realizzare applicazioni multipiattaforma senza richiedere hardware
proprietario.

\subsection{FastAPI}
FastAPI è un framework web open source per la realizzazione di API in Python.
È basato sulle annotazioni di tipo del linguaggio e utilizza Starlette per le
funzionalità web e Pydantic per la gestione dei dati. Supporta applicazioni
asincrone, API HTTP, WebSocket e generazione automatica della documentazione
OpenAPI.

\subsection{Uvicorn e ASGI}
Uvicorn è un server web per applicazioni Python conformi allo standard ASGI
(\textit{Asynchronous Server Gateway Interface}). ASGI definisce l'interfaccia
tra server e applicazione e, a differenza del precedente standard sincrono
WSGI, supporta connessioni di lunga durata e comunicazioni asincrone come i
WebSocket.

\subsection{WebSocket}
WebSocket è un protocollo di comunicazione che stabilisce un canale persistente
e bidirezionale tra client e server. Dopo una fase iniziale di handshake, le
due parti possono inviare messaggi in modo indipendente sulla stessa
connessione. Il protocollo è indicato con \texttt{WS}; la variante
\texttt{WSS} aggiunge la protezione crittografica fornita da TLS.

\subsection{JSON}
JSON (\textit{JavaScript Object Notation}) è un formato testuale per lo scambio
di dati strutturati. Rappresenta le informazioni mediante oggetti, array,
stringhe, numeri, valori booleani e valori nulli. È indipendente dal linguaggio
di programmazione ed è ampiamente utilizzato nelle API web.

\subsection{Base64}
Base64 è una codifica che rappresenta dati binari attraverso un insieme di
caratteri ASCII. Consente quindi di inserire immagini, audio o altri contenuti
binari in formati testuali come JSON. Non comprime né cifra le informazioni e
produce una rappresentazione più grande rispetto ai dati originali.

\subsection{HTTPX}
HTTPX è un client HTTP per Python dotato di interfacce sincrone e asincrone.
Supporta funzionalità quali connessioni persistenti, timeout, streaming delle
risposte e protocolli HTTP/1.1 e HTTP/2. È pensato per effettuare richieste di
rete senza rinunciare al modello asincrono del linguaggio.

\section{Validazione e osservabilità}

\subsection{Pydantic}
Pydantic è una libreria per la validazione, la serializzazione e il parsing dei
dati basata sulle annotazioni di tipo di Python. Attraverso modelli dichiarativi
consente di verificare la struttura dei valori in ingresso, convertirli nei
tipi attesi e produrre rappresentazioni compatibili con JSON. Può essere
utilizzata anche per definire configurazioni tipizzate e immutabili.

\subsection{Loguru}
Loguru è una libreria di logging per Python. Fornisce un'interfaccia compatta
per registrare eventi con differenti livelli di severità e permette di
configurare più destinazioni, dette \textit{sink}. Supporta inoltre rotazione,
conservazione e compressione automatica dei file di log.

\section{Elaborazione delle immagini}

\subsection{OpenCV}
OpenCV (\textit{Open Source Computer Vision Library}) è una libreria open
source dedicata alla computer vision e all'elaborazione delle immagini.
Comprende algoritmi per acquisizione, codifica, trasformazioni geometriche,
filtraggio, disegno, analisi di immagini e video. Offre interfacce per diversi
linguaggi, tra cui Python e C++.

\subsection{NumPy}
NumPy è la libreria fondamentale per il calcolo numerico in Python. Introduce
array multidimensionali omogenei e operazioni vettoriali eseguite in modo
efficiente. Immagini, coordinate e tensori numerici possono così essere
rappresentati ed elaborati senza ricorrere a cicli Python per ogni singolo
elemento.

\subsection{scikit-image e SSIM}
scikit-image è una libreria open source di algoritmi per l'elaborazione delle
immagini costruita sull'ecosistema scientifico di Python. Tra le metriche che
fornisce è presente SSIM (\textit{Structural Similarity Index Measure}), che
stima la somiglianza tra due immagini considerando luminanza, contrasto e
struttura locale.

\section{Object detection}

\subsection{RF-DETR}
RF-DETR (\textit{Roboflow Detection Transformer}) è un'architettura
transformer sviluppata da Roboflow per l'object detection in tempo reale. Usa
un backbone visuale DINOv2 per estrarre le caratteristiche dell'immagine e
restituisce, per ogni oggetto rilevato, una classe, una confidence e una
bounding box.

\subsection{PyTorch}
PyTorch è un framework open source per il machine learning e il deep learning.
Fornisce tensori multidimensionali, differenziazione automatica e strumenti per
definire, addestrare ed eseguire reti neurali. Le operazioni possono essere
eseguite su CPU oppure accelerate mediante GPU compatibili.

\subsection{Supervision}
Supervision è una libreria di utilità per applicazioni di computer vision.
Offre strutture dati e funzioni per rappresentare, filtrare, trasformare e
visualizzare i risultati prodotti dai modelli di detection, come bounding box,
classi e valori di confidence.

\section{Rilevamento della mano}

\subsection{MediaPipe Hand Landmarker}
MediaPipe è un framework multipiattaforma per la costruzione di pipeline di
machine learning applicate a dati multimediali. La componente Hand Landmarker
individua le mani presenti in un'immagine o in un flusso video e restituisce 21
punti anatomici per ciascuna mano, oltre alla lateralità. I risultati includono
sia coordinate normalizzate rispetto all'immagine sia coordinate
tridimensionali espresse nel sistema di riferimento della mano.

\section{Modelli visivo-linguistici}

\subsection{Vision Language Model}
Un Vision Language Model (VLM) è un modello multimodale capace di elaborare
congiuntamente informazioni visive e linguistiche.

\subsection{OpenRouter}
OpenRouter è una piattaforma di inferenza che offre un'unica API per accedere a
modelli generativi appartenenti a fornitori differenti. Espone un'interfaccia
compatibile con il formato delle API OpenAI e uniforma struttura delle
richieste, autenticazione e risposte. Può inoltre instradare le richieste tra
provider diversi e applicare meccanismi di fallback. OpenRouter non è quindi
un modello di intelligenza artificiale, ma un livello di accesso e
orchestrazione dei modelli.

\subsection{Gemini 2.5 Flash}
Gemini 2.5 Flash è un modello generativo multimodale sviluppato da Google. Può
ricevere in ingresso testo, immagini, audio e video e produrre testo. La
variante Flash è progettata per offrire un compromesso tra capacità del
modello, velocità di risposta e costo computazionale, risultando adatta a
carichi interattivi e ad alto volume.

\section{Sintesi vocale}

La sintesi vocale, o Text-to-Speech (TTS), è il processo che trasforma un testo
scritto in un segnale audio contenente parlato artificiale. I sistemi TTS
neurali apprendono pronuncia, prosodia, ritmo e caratteristiche della voce da
grandi quantità di registrazioni.

\subsection{ElevenLabs}
ElevenLabs è una piattaforma cloud dedicata alla generazione e all'elaborazione
della voce mediante modelli di intelligenza artificiale. La sua tecnologia
Text-to-Speech converte testo in parlato, supporta più lingue e consente di
selezionare voci con differenti caratteristiche espressive. Le funzionalità
sono accessibili tramite interfaccia web e API.

\subsection{Chatterbox TTS}
Chatterbox è una famiglia di modelli Text-to-Speech open source sviluppata da
Resemble AI. I modelli generano parlato a partire dal testo e supportano
funzionalità quali sintesi multilingue, controllo dell'espressività e
clonazione della voce. Essendo distribuito come software open source, può
essere eseguito su infrastruttura locale.

\section{Test automatici}

\subsection{pytest}
pytest è un framework per la scrittura e l'esecuzione di test in Python.
Utilizza normali funzioni e istruzioni \texttt{assert}, offre meccanismi di
scoperta automatica dei test e permette di condividere dati e risorse mediante
fixture.

\subsection{pytest-asyncio}
pytest-asyncio è un'estensione di pytest per il collaudo di codice basato su
\texttt{asyncio}. Fornisce un event loop controllato dall'ambiente di test e
consente di eseguire direttamente coroutine e fixture asincrone.
